\scanpage{325}
\margin{326}
3) Donnée d'une application (pour tt $p\geq0$)
\[
\operatorname{Hom}_h(A,M)\xrightarrow{\unsure{u_A^{(p)}|u_A^p}}\operatorname{Hom}_h(A,\Gamma^p M)
\]
(pas \inserted{néc.} lin) fonctorielle en $M$, notée
ce qui \deleted{fait} sera noté
\[
f\longmapsto f^{(p)}
\]
et satisfaisant les relations $f^{(0)}=1$, $f^{(1)}=f$, et
\[
\left\{\begin{aligned}
(f+g)^{(p)}&=\sum_{\unsure{i+j=0|i+j=p}} f^{(i)}g^{(j)}\\
(\lambda f)^{(p)}&=\lambda^p f^p\quad\text{(automatique!)}
\end{aligned}\right.
\]
et ---

4) Donnée d'une loi de p.d. sur
\[
\operatorname{Hom}_h(A,\Gamma^+M)\subset\operatorname{Hom}_h(A,M)
\]
$\forall M$, fonctorielle en $M$, et telle que

5) Loi à p.d. sur $\operatorname{Hom}(A,\Gamma^+A)\subset\operatorname{Hom}(A,\Gamma^*A)$
\margin{(\unsure{en. commut|an. commut})}

6) Donnée, pour tte $h$-algèbre $B$ (avec
idéal $J$ à p.d, d'une loi de
p.d. sur $\operatorname{Hom}_h(A,J)\subset\operatorname{Hom}_h(A,B)$,
fonctorielle en $(B,J)$ $\ldots\ldots$
